% Part: incompleteness
% Chapter: theories-computability
% Section: first-incompleteness

\documentclass[../../../include/open-logic-section]{subfiles}

\begin{document}

\olfileid{inc}{tcp}{inc}

\olsection{$\Th{Q}$ هیچ گسترش کامل و سازگاری ندارد که !!^{axiomatizable} باشد}

\begin{thm}
\ollabel{thm:first-incompleteness}
هیچ گسترش کامل و سازگاری از~$\Th{Q}$ وجود ندارد که
!!{axiomatizable} باشد.
\end{thm}

\begin{proof}
از پیش می‌دانیم که هیچ گسترش سازگار و تصمیم‌پذیری از~$\Th{Q}$ وجود
ندارد. اما اگر~$\Th{T}$ کامل و !!{axiomatizable} باشد، تصمیم‌پذیر است.
\end{proof}

\begin{explain}
این قضیه چندان از صورت‌بندی اصلی قضیهٔ نخست ناتمامیت گودل در سال ۱۹۳۱
دور نیست. گذشته از اصطلاحات جدیدتر، تفاوت‌های اصلی چنین‌اند: گودل به
جای «سازگار» از «$\omega$-سازگار» استفاده می‌کند؛ و نمی‌توانست با
کلیت کامل از «!!{axiomatizable}» سخن بگوید، زیرا مفهوم صوری
محاسبه‌پذیری هنوز شکل نگرفته بود. (مدل‌های صوری محاسبه‌پذیری در دههٔ
پس از آن، از جمله به دست خود گودل، و تا حد زیادی برای توصیف دقیق
نظریه‌هایی پدید آمدند که در معرض پدیدهٔ گودل قرار می‌گیرند.)

قضیه می‌گوید نمی‌توان هر سه ویژگیِ کامل‌بودن، سازگاری و
!!{axiomatizability} را هم‌زمان داشت. بااین‌حال، اگر هر یک از این سه
را کنار بگذارید، دو ویژگی دیگر دست‌یافتنی‌اند: $\Th{Q}$ سازگار و به‌طور
محاسبه‌پذیر اصل‌موضوعی‌شده است، اما کامل نیست؛ نظریهٔ ناسازگار کامل و
به‌طور محاسبه‌پذیر اصل‌موضوعی‌شده است (مثلاً با
$\{ \eq/[0][0] \}$)، اما سازگار نیست؛ و مجموعهٔ جمله‌های صادق حساب کامل
و سازگار است، اما به‌طور محاسبه‌پذیر اصل‌موضوعی‌شده نیست.
\end{explain}
\end{document}
